\section{Lösbarkeit von linearen Gleichungssystemen}
\label{sec:LoesbarkeitVonLGS}
In den Übungen zum vorangegangenen Kapitel sind wir auf die folgenden Gleichungssysteme gestossen:

\[
\begin{cases}
-4x + 6y = 5 \\
6x - 9y = -8
\end{cases}
\quad \text{und} \quad
\begin{cases}
-4x + 6y = 5 \\
6x - 9y = -7.5
\end{cases}
\]

Da 
\[
\det \begin{pmatrix} -4 & 6 \\ 6 & -9 \end{pmatrix} = 0,
\]
kann man diese Gleichungssysteme nicht mit der Inversen der Matrix lösen.

Wenn man sie auf konventionelle Art löst, stellt man leicht fest, dass das linke Gleichungssystem keine Lösung hat, das rechte hingegen unendlich viele.

\medskip

In diesem Kapitel werden wir Begriffe und Ideen einführen, mit denen wir einfach feststellen können, ob so ein System lösbar ist.

Zunächst einmal wollen wir die Situation geometrisch auffassen:  
Gesucht ist ein Vektor
\[
\vec{v} = \begin{pmatrix} x \\ y \end{pmatrix}
\]
mit der Eigenschaft, dass das Bild \( A\vec{v} \) unter der Matrix \( A \) gerade der Vektor
\[
\begin{pmatrix} 5 \\ -8 \end{pmatrix} \quad \text{bzw.} \quad \begin{pmatrix} 5 \\ -7.5 \end{pmatrix}
\]
ist.

Allgemein werden wir uns fragen, ob es zu einem Vektor \( \vec{b} \) einen Vektor \( \vec{v} \) gibt, so dass
\[
A\vec{v} = \vec{b}
\]

\subsection{Homogene Systeme}
\label{sec:HomogeneSysteme}
Wir beginnen mit dem einfachsten Fall, den sogenannten \emph{homogenen Systemen}:
\begin{definition}{Homogenes Gleichungssystem}{}
Ein \emph{homogenes lineares Gleichungssystem} ist ein System der Form
\[
A\vec{v} = \vec{0} \\
\]

\end{definition}

Wenn die Matrix $A$ invertierbar ist, dann hat das Gleichungssystem nur den Nullvektor als Lösung. 
In diesem Fall gilt nämlich:
$$ \vec{v} = A^{-1}\cdot \vec{0} = \vec{0}.$$

\begin{bsp}{}{}
Sei $A = \begin{pmatrix} 1 & 2 \\ 3 & 4 \end{pmatrix}$. Dann ist $A$ invertierbar ($\det A = -2 \neq0$), und das Gleichungssystem
\[
\begin{cases}
x+2y = 0 \\
3x+4y = 0
\end{cases}
\]
hat nur die triviale Lösung $\begin{pmatrix}
    x \\ y 
\end{pmatrix} = \begin{pmatrix} 0 \\ 0 \end{pmatrix}$.
\end{bsp}

Wenn die Matrix $A$ nicht invertierbar ist, dann hat das Gleichungssystem unendlich viele Lösungen. Diese Tatsache soll in der folgenden Lernaufgabe untersucht werden.

\begin{LA}{Homogene LGS mit singulärer Matrix}{}
    Betrachten Sie zuerst eine konkrete Matrix $A$:
    \[
    A = \begin{pmatrix} 1 & 2 \\ 2 & 4 \end{pmatrix}
    \]
    \begin{enumerate}
        \item Zeigen Sie, dass $A$ nicht invertierbar ist.
        \item Zeigen Sie, dass der Vektor $\vec{v} = \begin{pmatrix}
            2 \\ -1
        \end{pmatrix}$ eine Lösung des Gleichungssystems
        $
        A\vec{v} = \vec{0}
        $
        ist.
        \item Zeigen Sie, dass auch der Vektor $\vec{v} = \begin{pmatrix}
            4 \\ -2
        \end{pmatrix}$ eine Lösung des Gleichungssystems
        $
        A\vec{v} = \vec{0}
        $
        ist.
        \item Begründen Sie, dass es unendlich viele Lösungen des Gleichungssystems
        $
        A\vec{v} = \vec{0}
        $
        gibt.
    \end{enumerate}
    Als nächstes untersuchen wir die allgemeine Situation. Sei $A$ die Matrix
    $$ A = \begin{pmatrix} a & b \\ c & d \end{pmatrix}$$
    mit $\det A = ad - bc = 0$.
    \begin{enumerate}[resume]
        \item Zeigen Sie dass der Vektor $\vec{v} = \begin{pmatrix}
            -b \\ a
        \end{pmatrix}$ eine Lösung des Gleichungssystems
        $
        A\vec{v} = \vec{0}
        $
        ist.
        \item Zeigen Sie, dass auch der Vektor $\vec{v} = \begin{pmatrix}
            d \\ -c
        \end{pmatrix}$ eine Lösung des Gleichungssystems
        $
        A\vec{v} = \vec{0}
        $
        ist.
        \item Begründen Sie, dass es unendlich viele Lösungen des Gleichungssystems
        $
        A\vec{v} = \vec{0}
        $
        gibt.
        \begin{enumerate}
            \item Für den Fall, dass $a=b=c=d=0$ ist.
            \item Für den Fall, dass ein Eintrag von $A$ ungleich null ist.
                \end{enumerate}
    \end{enumerate}
\end{LA}
\clearpage
\begin{definition}
    {Kern}{}

Sei \( A \) eine Matrix.  

Die Menge aller Vektoren, die auf den Nullvektor abgebildet werden, nennt man den \emph{Kern} von \( A \) und schreibt dafür:
\[
\ker(A) = \left\{ \vec{v} \in \mathbb{R}^2 \mid A\cdot\vec{v} = \vec{0} \right\}
\]
Der Kern ist die Menge aller Lösungen des homogenen Gleichungssystems
\[
A\cdot\vec{v} = \vec{0} = \vec{0}
\]
\end{definition}
\begin{satz}
    {}{}
    Für den Kern einer $2\times 2$-Matrix $A$ gilt:
    \begin{enumerate}
        \item Falls $\det A \neq 0$ ist $\ker (A) = \{0\}$
        \item Falls $A = \begin{pmatrix}
            0 & 0 \\
            0 & 0 \end{pmatrix}$ ist $\ker (A) = \mathbb{R}^2$
        \item In allen anderen Fällen ist $\ker (A)$ eine Gerade durch den Ursprung.
    \end{enumerate}
\end{satz}

\begin{bsp}
{Kern von $B = \begin{pmatrix} 1 & 3 \\ 4 & -1 \end{pmatrix}.$}{}


Es gilt $\det B = -13 \neq 0$, also 

\[
\ker(B) = \left\{ \vec{0} \right\}.
\]

\end{bsp}


\begin{bsp}
    {Kern von $A = \begin{pmatrix} 1 & 3 \\ 3 & 9 \end{pmatrix}$ }{}


Es gilt $\det A = 0$, also ist $A$ nicht invertierbar.

Dreimal die erste Spalte minus die zweite ergibt den Nullvektor, also ist:
\[
\ker(A) = \left\{ t \begin{pmatrix} 3 \\ -1 \end{pmatrix} \;\middle|\; t \in \mathbb{R} \right\}.
\]

\end{bsp}


\subsection{Inhomogene Systeme}
\label{sec:InhomogeneSysteme}
\begin{definition}{Inhomogenes Gleichungssystem}{}
Ein \emph{inhomogenes lineares Gleichungssystem} ist ein System der Form
\[
A\vec{v} = \vec{b} \quad \text{mit } \vec{b} \neq \vec{0}
\]
\end{definition}
Wenn die Matrix $A$ invertierbar ist, dann hat das Gleichungssystem genau eine Lösung:
\[
\vec{v} = A^{-1}\cdot \vec{b}.\]
\begin{bsp}
{Inhomogenes LGS mit invertierbarer Matrix}{}
Sei $A = \begin{pmatrix} 1 & 2 \\ 3 & 4 \end{pmatrix}$ und $\vec{b} = \begin{pmatrix} 5 \\ 6 \end{pmatrix}$. Dann ist $A$ invertierbar ($\det A = -2 \neq0$), und das Gleichungssystem
\[
\begin{cases}
x+2y = 5 \\
3x+4y = 6
\end{cases}
\]
hat genau eine Lösung $$\begin{pmatrix}
    x \\ y
\end{pmatrix} = A^{-1} \cdot \begin{pmatrix}
    5 \\ 6
\end{pmatrix} = \frac{1}{-2} \begin{pmatrix}
    4 & -2 \\ -3 & 1
\end{pmatrix}\cdot \begin{pmatrix}
    5 \\ 6
\end{pmatrix} = \begin{pmatrix} -4 \\ 4.5 \end{pmatrix}$$
\end{bsp}

Wenn die Matrix $A$ nicht invertierbar ist, dann hat das Gleichungssystem entweder keine Lösung oder unendlich viele Lösungen. Diese Tatsache soll wieder in einer Lernaufgabe untersucht werden.

